\epanalipsi{1}

\tmima{Λεξικό του κεφαλαίου}

\renewcommand{\arraystretch}{1.5}
\begin{tabular*}{\linewidth}{@{\extracolsep{\fill}} >{\raggedright\arraybackslash}p{0.26\linewidth} >{\raggedright\arraybackslash}p{0.64\linewidth}}
\toprule
\textbf{Λέξη ή τύπος} & \textbf{Με μία πρόταση}\\
\midrule
φυσικός & Ένας από τους $0,\,1,2,\,3,\dots$: μετράει πράγματα που δεν σπάνε.\\
στρογγυλοποίηση & Αλλάζω έναν αριθμό με κοντινό του, με συμφωνημένα βήματα, μία φορά.\\
$\alpha^{\nu}$ & Ο $\alpha$ γραμμένος $\nu$ φορές σε γινόμενο.\\
προτεραιότητα & Παρενθέσεις, δυνάμεις, $\cdot$ και $:$, και τέλος $+$ και $-$.\\
$\Delta=\delta\cdot\pi+\upsilon$ & Ευκλείδεια διαίρεση, με $\upsilon<\delta$.\\
κριτήρια & Κανόνες στα ψηφία που δείχνουν αν ένας αριθμός διαιρείται με το 2, 3, 4, 5, 9, 10, 25.\\
πρώτος & Ο αριθμός μεγαλύτερος από το 1 με διαιρέτες μόνο τον 1 και τον εαυτό του.\\
ΜΚΔ & Ο μεγαλύτερος κοινός διαιρέτης: μοιράζω.\\
ΕΚΠ & Το μικρότερο (μη μηδενικό) κοινό πολλαπλάσιο: συναντιέμαι.\\
\bottomrule
\end{tabular*}

\tmima{Ασκήσεις επανάληψης}

\ex Γράψε με ψηφία τον αριθμό «τριάντα επτά εκατομμύρια, δύο χιλιάδες, πέντε». Μετά στρογγυλοποίησέ τον στα εκατομμύρια.

\ex Βάλε σε αύξουσα σειρά τους 5.012, 5.102, 5.021, 4.999 και 10.000.

\ex Στρογγυλοποίησε το 86.549 στις δεκάδες, στις εκατοντάδες και στις χιλιάδες.

\ex Υπολόγισε με τις ιδιότητες των πράξεων: α) $25\cdot 17\cdot 4$ \quad β) $99\cdot 8$ \quad γ) $36\cdot 15+36\cdot 85$ \quad δ) $47+68+53+32$.

\ex Βρες την τιμή: α) $3\cdot(8-2)^2-10$ \quad β) $2^5-5^2$ \quad γ) $(2+3)\cdot 4-2^3$ \quad δ) $100-4\cdot 5^2$.

\ex Κάνε την Ευκλείδεια διαίρεση $1.000:48$ και γράψε την ισότητα $\Delta=\delta\cdot\pi+\upsilon$. Ποιος αριθμός δίνει πηλίκο 7 και υπόλοιπο 6 όταν διαιρεθεί με το 15;

\ex Σήμερα είναι Πέμπτη. Τι μέρα θα είναι μετά από 100 ημέρες;

\ex Με ποιους από τους 2, 3, 4, 5, 9, 10 και 25 διαιρείται ο αριθμός 3.240;

\ex Ποιο ψηφίο πρέπει να μπει στη θέση του $\square$ ώστε ο $5\square 7$ να διαιρείται με το 9;

\ex Ποιοι από τους 57, 59, 63 και 67 είναι πρώτοι; Για τους σύνθετους, βρες την ανάλυσή τους.

\ex Ανάλυσε τον 360 σε γινόμενο πρώτων παραγόντων. Πόσους πρώτους διαφορετικούς παράγοντες έχει;

\ex Βρες το ΜΚΔ και το ΕΚΠ των αριθμών 36 και 60.

\ex \textbf{Οι κάρτες του Στέλιου.} Ο Στέλιος έχει 36 κάρτες με ζώα και 60 κάρτες με αυτοκίνητα. Θέλει να φτιάξει όσο το δυνατόν περισσότερες ίδιες στοίβες, με και από τα δύο είδη, χωρίς να περισσέψει καμία. Πόσες στοίβες θα γίνουν και τι έχει η κάθε μία;

\ex \textbf{Η πισίνα.} Ο Ηλίας πάει στην πισίνα κάθε 8 μέρες και η Ζωή κάθε 12. Σήμερα πήγαν μαζί. Πότε θα ξαναπάνε μαζί;

\ex Σωστό ή λάθος; α) Κάθε άρτιος αριθμός διαιρείται με το 4. \quad β) $\MKD(12,\,18)=6$. \quad γ) Το 91 είναι πρώτος. \quad δ) $2^4=8$. \quad ε) $17-(9-3)=17-9-3$. \quad στ) $0:5=0$.

\ex Ο Νίκος λέει ότι $6+4\cdot 5=50$ και η Ζωή ότι $6+4\cdot 5=26$. Ποιος έχει δίκιο και γιατί; Πού πρέπει να βάλεις παρένθεση ώστε να βγει 50;

\ex Πόσους διαιρέτες έχει το 100; Πόσοι από αυτούς είναι πρώτοι;

\ex \textbf{Ερωτήσεις αυτοαξιολόγησης (πράξεις και δυνάμεις).} Σωστό ή λάθος; α) $(50-20)-5=50-(20-5)$. \quad β) $5\cdot 5\cdot 5=5^3$. \quad γ) $3^4=12$. \quad δ) $\beta+\beta+\beta=3\beta$. \quad ε) $2^4+4=20$. \quad στ) $2\cdot 10^2+5\cdot 10+1=251$. \quad ζ) $30-18:6=2$. \quad η) $6\cdot 4-3+7=28$.

\ex \textbf{Ερωτήσεις αυτοαξιολόγησης (διαιρετότητα).} Σωστό ή λάθος; α) Το άθροισμα δύο περιττών αριθμών είναι πάντα άρτιος. \quad β) Αν ο $\alpha$ είναι πολλαπλάσιο του $\beta$ (με $\beta\neq 0$), τότε ο $\beta$ διαιρεί τον $\alpha$. \quad γ) Το 51 είναι πολλαπλάσιο του 3 και του 5. \quad δ) Ο 630 διαιρείται με το 9. \quad ε) Η διαίρεση $365:20$ δίνει πηλίκο 18 και υπόλοιπο 5. \quad στ) Η σχέση $250=12\cdot 20+10$ είναι Ευκλείδεια διαίρεση. \quad ζ) Ο αριθμός $4\alpha+8$ διαιρείται με το 4. \quad η) Ο 312 διαιρείται με το 4 και με το 8.

\ex \textbf{Ερωτήσεις αυτοαξιολόγησης (ΜΚΔ, ΕΚΠ, πρώτοι).} Σωστό ή λάθος; α) Ο 14 και ο 42 έχουν ΜΚΔ το 7. \quad β) $\EKP(5,\,15)=75$. \quad γ) Η γραφή $150=5\cdot 30$ είναι ανάλυση σε γινόμενο πρώτων παραγόντων. \quad δ) Δύο διαδοχικοί φυσικοί (π.χ. 14 και 15) είναι πάντα πρώτοι μεταξύ τους. \quad ε) Αν ο $\alpha$ διαιρεί τον $\beta$, τότε $\MKD(\alpha,\,\beta)=\alpha$. \quad στ) Ο 2 είναι ο μόνος άρτιος πρώτος.

\tmima{Ο χάρτης της επανάληψης}

Μετά τις λύσεις, δες σε ποιες ασκήσεις δυσκολεύτηκες. Ο πίνακας σου λέει ποιο κεφάλαιο να ξαναδιαβάσεις.

\renewcommand{\arraystretch}{1.5}
\begin{tabular*}{\linewidth}{@{\extracolsep{\fill}} >{\raggedright\arraybackslash}p{0.38\linewidth} >{\raggedright\arraybackslash}p{0.52\linewidth}}
\toprule
\textbf{Αν δυσκολεύτηκες στις} & \textbf{ξαναδιάβασε τις παραγράφους}\\
\midrule
Ε1.1, Ε1.2, Ε1.3 & 1.1 (φυσικοί, διάταξη, στρογγυλοποίηση)\\
Ε1.4 & 1.2 (ιδιότητες των πράξεων)\\
Ε1.5, Ε1.16, Ε1.18 & 1.3 (δυνάμεις, προτεραιότητα)\\
Ε1.6, Ε1.7, Ε1.19 & 1.4 (Ευκλείδεια διαίρεση)\\
Ε1.8, Ε1.9 & 1.5, μέρος Α (κριτήρια διαιρετότητας)\\
Ε1.10, Ε1.11, Ε1.17 & 1.5, μέρος Β (πρώτοι, ανάλυση)\\
Ε1.12, Ε1.13, Ε1.14 & 1.5, μέρος Γ (ΜΚΔ, ΕΚΠ)\\
Ε1.15, Ε1.18, Ε1.19, Ε1.20 & 1.1 έως 1.5\\
\bottomrule
\end{tabular*}
